\documentclass[preprint,12pt]{elsarticle}

\usepackage{amsmath,amssymb,amsfonts}
\usepackage{booktabs}
\usepackage{hyperref}
\usepackage{url}
\usepackage{listings}
\usepackage{xcolor}

\hypersetup{
    colorlinks=true,
    linkcolor=blue,
    citecolor=blue,
    urlcolor=blue
}

\journal{SoftwareX}

\begin{document}

\begin{frontmatter}

\title{GeoSolver 2D: A Provenance-Aware Dynamic Geometry Engine with Dual-Regime and Permutation-Invariant Classification}

\author{Bao-Trung Nguyen\corref{cor1}}
\ead{trungyds16@gmail.com}
\cortext[cor1]{Corresponding author.}

\affiliation{organization={Independent Researcher},
            city={Ho Chi Minh City},
            country={Vietnam}}

\begin{abstract}
Interactive Dynamic Geometry Systems (DGS) are prone to geometric misclassifications when users construct or perturb configurations using freehand gestures. Floating-point rounding errors and non-transitive relative tolerances frequently lead to false-positive recognitions (such as misclassifying near-regular integer-lattice triangles as equilateral) or permutation-dependent outcomes (where the classified type depends on vertex enumeration order). We introduce \textbf{GeoSolver 2D}, an open-source geometry engine designed to remove these numerical artifacts in triangle classification. GeoSolver 2D implements: (1) a \textit{Dynamic Taint Tracking Engine} that enforces an acyclic dependency graph across interactive constructions using an iterative topological traversal with explicit stack safety and strictly linear $O(V+E)$ complexity; (2) a dual-regime validation model distinguishing between exact integer-lattice coordinates ($\mathbb{Z}^2$) and continuous approximations; and (3) a \textit{Specificity-Ordered Classifier} operating on scale-normalized edge lengths and cosine predicates, which guarantees bitwise vertex-permutation invariance for type classification and equivariance for apex labeling. Scale-invariance holds unconditionally within the continuous regime, with explicit provenance tagging preventing artificial scale-induced transitions. Evaluated across an automated verification suite of 62 regression tests (27 headless core library tests and 35 desktop integration tests) and an empirical benchmark suite achieving 91.4\% accuracy on canonical prototypes, the engine demonstrates bitwise permutation invariance, strict apex equivariance, and bounded linear $O(V+E)$ traversal on cycle-free construction graphs. The core computation engine is released as open-source software under the GNU General Public License v3 (GPLv3).
\end{abstract}

\begin{keyword}
Dynamic Geometry System \sep Taint Tracking \sep Permutation Invariance \sep Integer Lattice \sep Geometric Classification \sep Software Verification
\end{keyword}

\end{frontmatter}

\section*{Metadata Table}

\begin{table}[h!]
\centering
\small
\begin{tabular}{lp{9.5cm}}
\toprule
\textbf{Current code version} & v3.2.1 \\
\midrule
\textbf{Permanent code repository} & \url{https://github.com/trungyds16-coder/geosolver-core} \\
\textbf{Legal Code License} & GNU General Public License v3 (GPLv3) \\
\textbf{Code versioning system used} & Git (Release tag: \texttt{v3.2.1}) \\
\textbf{Software code languages} & Python 3.10+ \\
\textbf{Compilation requirements} & None (Pure Python standard library) \\
\textbf{Supported operating systems} & Windows 10/11, Linux (Ubuntu 22.04+), macOS; CI tested \\
\textbf{Automated test suite} & Python \texttt{unittest} (27 headless CI tests, 35 GUI tests, 100\% PASS) \\
\textbf{Persistent identifier (DOI)} & \url{https://doi.org/10.5281/zenodo.24249595} \\
\textbf{Documentation \& Support} & \url{https://github.com/trungyds16-coder/geosolver-core} \\
\bottomrule
\end{tabular}
\caption{Software metadata for GeoSolver 2D.}
\label{tab:metadata}
\end{table}

\section{Motivation and Significance}

\subsection{Background and Problem Statement}
Dynamic Geometry Systems (DGS) are foundational tools in computer-assisted mathematics education, automated theorem proving, and computer-aided design (CAD). However, interactive geometry environments operate under a fundamental tension: \textbf{the gap between continuous user gestures (such as freehand mouse perturbations) and exact discrete mathematical truth}.

When users manipulate points on an interactive canvas, coordinate values undergo floating-point perturbations ($x \in \mathbb{R}$). Conventional DGS implementations attempt to identify geometric configurations (such as right, isosceles, or equilateral triangles) by evaluating heuristic tolerance thresholds ($\varepsilon$). In practice, this design introduces two major failure modes:
\begin{enumerate}
    \item \textbf{Misclassification on Integer Lattices ($\mathbb{Z}^2$):} By the classical Shoelace Formula (Gauss Area Formula), any triangle whose vertices lie on the integer lattice $\mathbb{Z}^2$ has a strictly rational area:
    \begin{equation}
    \text{Area} = \frac{1}{2} |X_A(Y_B - Y_C) + X_B(Y_C - Y_A) + X_C(Y_A - Y_B)| \in \frac{1}{2}\mathbb{Z} \subset \mathbb{Q}
    \end{equation}
    Conversely, an equilateral triangle with side length $a$ requires an area of $\text{Area} = \frac{\sqrt{3}}{4}a^2$. Because the squared distance between any two integer vertices is an integer ($a^2 = d_{\mathbb{Z}}^2 \in \mathbb{Z}$), the area of a lattice equilateral triangle would be strictly irrational due to $\sqrt{3} \notin \mathbb{Q}$. Since an irrational number cannot equal a rational number, \textbf{the equilateral hypothesis is provably unreachable on $\mathbb{Z}^2$}. Heuristic DGS engines that blindly apply continuous relative tolerances to lattice points can misclassify near-equilateral grid triangles as strictly equilateral, which may give students incorrect mathematical feedback.
    \item \textbf{Permutation Dependence from Non-Transitive Tolerances:} In metric spaces over the real numbers, relative tolerance relations $|a - b| \le \varepsilon$ do not satisfy transitivity:
    \begin{equation}
    |a - b| \le \varepsilon \land |b - c| \le \varepsilon \not\implies |a - c| \le \varepsilon
    \end{equation}
    In floating-point arithmetic compliant with IEEE-754~\cite{ieee754}, when an interactive engine evaluates candidate hypotheses through sequential conditional logic (\texttt{if/elif}), the resulting classification depends on the arbitrary indexing order of the vertices ($\sigma(A, B, C)$ vs $\sigma(B, C, A)$).
\end{enumerate}

\subsection{Related Work \& Positioning}
In computational geometry, the standard defense against numerical instability is \textbf{Exact Geometric Computation (EGC)}, exemplified by Shewchuk's adaptive-precision floating-point predicates~\cite{shewchuk1997} and the CGAL geometric kernel~\cite{cgal2024}. EGC solves orientation and incircle tests robustly, but requires rigorous floating-point filtering with multi-precision fallback.

To preserve computational simplicity while eliminating classification artifacts, GeoSolver 2D adapts concepts from dynamic taint analysis~\cite{schwartz2010} into geometric provenance tracking. By classifying objects based on their operational history rather than floating-point coordinates alone, GeoSolver 2D bridges the gap between interactive usability and formal mathematical truth.

\section{Software Architecture and Mathematical Formulation}

GeoSolver 2D partitions geometric processing into three decoupled layers:
\begin{enumerate}
    \item \textbf{Dynamic Taint Tracking Engine (\texttt{provenance.py}):} Evaluates derivation history across interactive Directed Acyclic Graphs (DAGs).
    \item \textbf{Dual-Regime Coordinate Model (\texttt{as\_scaled\_ints}):} Partitions inputs between discrete integer lattices ($\mathbb{Z}^2$) and continuous floating-point domains.
    \item \textbf{Specificity-Ordered Classifier (\texttt{classifier.py}):} Implements an invariant hypothesis sorting hierarchy with scale-normalized edge metrics and cosine predicates.
\end{enumerate}

\subsection{Dynamic Taint Tracking Engine}
Every geometric entity $E$ maintains a provenance flag:
\begin{equation}
\text{exact}(E) = \left(\bigwedge_{P \in \text{Parents}(E)} \text{exact}(P)\right) \land \neg \text{is\_tainted}(E)
\end{equation}
\begin{itemize}
    \item \textbf{Taint Conservation Law:} Freehand canvas gestures unconditionally inject taint (\texttt{exact=False}), which strictly propagates to all descendant constructions (midpoints, intersections, perpendicular feet). Snapping to the coordinate grid (\texttt{snap\_to\_grid}) restores exact integer status (\texttt{exact=True}, \texttt{source='grid\_snap'}). Stored flags on derived entities are never trusted over dynamic parent state.
    \item \textbf{DAG Cycle Safety \& Stack Safety:} Traversal is protected against cyclic dependencies via an active \texttt{\_visiting} tracking set. Traversal is evaluated iteratively via an explicit stack (topological post-order DFS), eliminating recursion depth limits and guaranteeing strict traversal-order invariance across subtrees.
\end{itemize}

\subsection{Dual-Regime Coordinate Model \& \texttt{as\_scaled\_ints}}
Input coordinates are evaluated via \texttt{as\_scaled\_ints} across scale factors $10^k$ ($k \in \{0, 1, 2, 3, 4\}$):
\begin{itemize}
    \item If all coordinates satisfy $|v \cdot 10^k - \text{round}(v \cdot 10^k)| < 10^{-5}$, they are converted to scaled integers $[\text{round}(v \cdot 10^k)]$ and evaluated under exact integer arithmetic on $\mathbb{Z}^2$.
    \item Untagged coordinate tuples default to exact if they align with the scale grid. Coordinates with freehand provenance unconditionally bypass integer scaling and are quarantined into the continuous floating-point regime until explicitly snapped to grid.
\end{itemize}

\subsection{Specificity Hierarchy and Invariance Theorem}
To prevent hypothesis suppression, GeoSolver 2D evaluates hypotheses through a strict specificity hierarchy:
\begin{equation}
\text{Equilateral} \succ \text{Right-Isosceles} \succ \{\text{Isosceles}, \text{Right}\} \succ \text{Scalene}
\end{equation}
Right angles are evaluated via the scale-correct cosine predicate:
\begin{equation}
|\cos \theta| = \left| \frac{\mathbf{u} \cdot \mathbf{v}}{\|\mathbf{u}\| \|\mathbf{v}\|} \right| \le 0.015 \quad (\approx |\theta - 90^\circ| \le 0.86^\circ)
\end{equation}

\textbf{Theorem 1 (Vertex-Permutation Invariance \& Apex Equivariance):}  
\textit{Let $V = (v_1, v_2, v_3)$ be a triangle with edge lengths $l_{12}, l_{23}, l_{31}$. For each vertex $v \in V$, define its local geometric record $R(v) = (\Delta(v), d^2_{\text{opp}}(v))$ where $\Delta(v)$ is the relative mismatch between the two sides incident to $v$, and $d^2_{\text{opp}}(v)$ is the squared length of the opposite side. Because IEEE-754 subtraction and multiplication satisfy $\text{fl}(a - b)^2 = \text{fl}(b - a)^2$ under round-to-nearest mode, each side length is bitwise symmetric. For any permutation $\sigma \in S_3$, the set of records $\{R(v) \mid v \in V\}$ is mapped bijectively to $\{R(\sigma(v))\}$. Because candidate hypotheses are sorted over the label-free key $(\Delta, d^2_{\text{opp}})$, the classified triangle type is strictly invariant under $\sigma$, and the apex selection satisfies $\text{Apex}(\sigma(V)) \equiv \sigma(\text{Apex}(V))$. In the discrete regime, two apex candidates can tie on opposite squared length $d^2_{\text{opp}}$ if and only if all three sides are equal ($a^2 = b^2 = c^2$); because equilateral configurations are unreachable on $\mathbb{Z}^2$ (Shoelace area formula), candidate keys are strictly unique. In the continuous regime, two non-winning candidates can tie on the candidate key if and only if their opposite sides are bitwise equal; however, the winning apex candidate is strictly unique because the third vertex then has leg mismatch $\Delta = 0$ and strictly wins the candidate sort. Equilateral configurations are resolved at the highest specificity level without apex labeling.}

\section{Illustrative Examples}
\begin{itemize}
    \item \textbf{Example 1 (Near-Equilateral Lattice Triangle on $\mathbb{Z}^2$):} Consider $A(0, 0)$, $B(1000, 0)$, and $C(500, 866)$. Pairwise squared distances: $d^2(AB) = 1,000,000$ and $d^2(AC) = d^2(BC) = 500^2 + 866^2 = 999,956$. Relative difference: $\frac{1000 - \sqrt{999,956}}{1000} \approx 0.0022\% \ll 0.6\%$. A tolerance-based continuous classifier would evaluate this as equilateral. GeoSolver 2D checks integer inequality $1,000,000 \ne 999,956$, provably rejects equilateral via the Shoelace area theorem, and correctly outputs \textbf{Isosceles with apex at $C$} across all 6 input permutations.
    \item \textbf{Example 2 (Permutation Equivariance on Near-Tie Isosceles Configuration):} Consider $P(0, 0)$, $Q(1, 0)$, and $R(0.51, 0.86598)$ where $PQ = 1.0$, $PR \approx 1.0050$, and $QR \approx 0.9950$. Both pairs $(PQ, PR)$ and $(PQ, QR)$ satisfy the $0.6\%$ tolerance. By sorting candidates via the label-free key $(\text{diff\_leg}, d^2_{\text{opposite}})$, GeoSolver 2D deterministically selects apex $P$ across all 6 input orderings, preserving physical equivariance.
\end{itemize}

\section{Empirical Verification \& Sensitivity Analysis}

\subsection{Canonical Benchmark, Parameter Sensitivity \& Stress Tests}
We evaluated tolerance sensitivity across 1,000 synthetic triangles (250 each of ground-truth Equilateral, Isosceles, Right, and Scalene prototypes) using \texttt{scripts/sensitivity\_benchmark.py} (\texttt{seed = 42}) perturbed by Gaussian coordinate jitter ($\sigma = 0.03$, 0.3\% of nominal edge length 10.0). Table~\ref{tab:benchmark} reports the classification metrics.

\begin{table}[h!]
\centering
\small
\begin{tabular}{lcccccc}
\toprule
\textbf{Triangle Class} & \textbf{Ground Truth} & \textbf{TP} & \textbf{Recall} & \textbf{Precision} & \textbf{F1-Score} \\
\midrule
Equilateral & 250 & 205 & 82.0\% & 100.0\% & 90.1\% \\
Isosceles & 250 & 212 & 84.8\% & 82.5\% & 83.6\% \\
Right & 250 & 247 & 98.8\% & 100.0\% & 99.4\% \\
Scalene & 250 & 250 & 100.0\% & 85.9\% & 92.4\% \\
\midrule
\textbf{Overall Dataset} & \textbf{1,000} & \textbf{914} & \textbf{91.4\%} & \textbf{91.4\%} & \textbf{91.4\%} \\
\bottomrule
\end{tabular}
\caption{Classification benchmark on 1,000 synthetic triangles ($\text{tol\_rel} = 0.006$, seed 42).}
\label{tab:benchmark}
\end{table}

\textit{Empirical sensitivity scan across relative edge tolerance $\text{tol\_rel} \in [0.002, 0.010]$:}
\begin{itemize}
    \item $\text{tol\_rel} = 0.002$: Overall = 62.9\%, Equilateral Recall = 15.2\%, Precision = 100.0\%, Scalene Precision = 55.7\%.
    \item $\text{tol\_rel} = 0.004$: Overall = 77.8\%, Equilateral Recall = 47.6\%, Precision = 100.0\%, Scalene Precision = 72.5\%.
    \item $\text{tol\_rel} = 0.006$ (Default): Overall = 91.4\%, Equilateral Recall = 82.0\%, Precision = 100.0\%, Scalene Precision = 85.9\%.
    \item $\text{tol\_rel} = 0.008$: Overall = 96.8\%, Equilateral Recall = 92.8\%, Precision = 100.0\%, Scalene Precision = 94.7\%.
    \item $\text{tol\_rel} = 0.010$: Overall = 98.5\%, Equilateral Recall = 97.6\%, Precision = 100.0\%, Scalene Precision = 96.5\%.
\end{itemize}

\textit{Hard Negatives Stress Test (Near-Boundary Asymmetric Triangles):} Table~\ref{tab:hard_negatives} evaluates 200 near-boundary scalene triangles (100 near-isosceles with side mismatch in $0.8\%–1.5\%$, and 100 near-right with angles in $[88.0^\circ, 88.8^\circ]$ where $|\cos \theta| \in [0.021, 0.035]$).

\begin{table}[h!]
\centering
\small
\begin{tabular}{lccccc}
\toprule
$\text{tol\_rel}$ & \textbf{Near-Iso FP} & \textbf{Near-Right FP} & \textbf{Total TN} & \textbf{Total FP} & \textbf{FP Rate} \\
\midrule
0.002 & 0/100 & 0/100 & 200/200 & 0/200 & 0.0\% \\
0.004 & 0/100 & 0/100 & 200/200 & 0/200 & 0.0\% \\
0.006 (Default) & 0/100 & 0/100 & 200/200 & 0/200 & 0.0\% \\
0.008 & 1/100 & 0/100 & 199/200 & 1/200 & 0.5\% \\
0.010 & 36/100 & 0/100 & 164/200 & 36/200 & 18.0\% \\
\bottomrule
\end{tabular}
\caption{Decoupled hard negatives stress test across tolerance thresholds.}
\label{tab:hard_negatives}
\end{table}

\textit{Comparative Baselines:}
\begin{enumerate}
    \item \textbf{Permutation Equivariance:} Across 200 near-tie configurations under all 6 permutations, a naive positional classifier exhibited 100.0\% apex-variant outcomes (200/200), whereas GeoSolver 2D achieved 0.0\% variance (0/200).
    \item \textbf{Lattice Exactness on $\mathbb{Z}^2$:} A continuous classifier suffered a 75.0\% false-positive equilateral rate (3/4). GeoSolver 2D proved non-existence with 0.0\% false positives (0/4).
    \item \textbf{Decimal Lattice Sweep ($k \in \{1, 2, 3, 4\}$):} Across 496 non-collinear candidate pairs, GeoSolver 2D achieved 100.0\% exact recognition (496/496: 490 isosceles + 6 right-isosceles, all exact: True with apex at $A$).
\end{enumerate}

\subsection{Automated Regression Test Suite}
\begin{itemize}
    \item \textbf{Headless Core Library Suite (27 tests):} Pure Python standard library, executed on GitHub Actions CI across Ubuntu, Windows, and macOS (Python 3.10–3.13). Covers integer arithmetic, Shoelace assertions, permutation invariance, DAG taint conservation, cycle safety, deep DAG order-independence, stack safety up to 2,000 nodes, and 25-level diamond ladder traversal with strict $O(V+E)$ complexity (52 evaluated nodes). Execution time: \textbf{0.019s} (PASS).
    \item \textbf{Desktop GUI Integration Suite (35 tests):} Covers Tkinter canvas events, mouse drag snapping, undo/redo instance rebinding, and WCAG 2.1 AA luminance contrast. Execution time: \textbf{10.826s} (PASS).
    \item \textbf{Total Test Execution:} 62 / 62 Tests PASS (100.0\%).
\end{itemize}

\section{Impact and Practical Applications}

GeoSolver 2D addresses a specific, well-defined gap: interactive geometry tools that classify freehand or lattice-based triangles using tolerance heuristics can return results that depend on vertex order or that contradict exact arithmetic. By making classification permutation-invariant and by evaluating integer-lattice inputs exactly, the engine is designed to give consistent, mathematically defensible feedback in two settings:

\begin{enumerate}
    \item \textbf{Educational Geometry Software:} Tolerance-based recognition can label an imperfect drawing as equilateral or right-angled. GeoSolver 2D is intended to reduce such misleading feedback by separating exact lattice inputs from freehand inputs and by reporting near-type deviations descriptively. Its effect on learning has not yet been evaluated.
    \item \textbf{Geometric Modelling Pipelines:} Because the classifier is invariant to vertex ordering, its output does not depend on import or serialization order, which simplifies testing and reproducibility in tools that infer constraints from point sets.
    \item \textbf{Current Use and Limitations:} The educational edition has been distributed through the Microsoft Store. No classroom study has been conducted; this release focuses on the correctness and reproducibility of the core engine. Evaluating the effect on students' reasoning, and usability with teachers, is planned future work. The core library is a dependency-free Python package under GPLv3, so it can be reused and extended independently of the desktop application.
\end{enumerate}

\section{Conclusions and Future Work}
GeoSolver 2D demonstrates that combining dynamic taint tracking with specificity-ordered predicates addresses the tension between continuous user interaction and discrete mathematical truth. By evaluating exact integer lattice arithmetic when appropriate and sorting candidate hypotheses deterministically under continuous perturbations, the engine guarantees intra-regime scale invariance and permutation-invariant classification. The present work validates the classification engine through unit tests and synthetic benchmarks; classroom evaluation remains future work. Future developments will also explore extending this architecture to 3D spatial geometry and polyhedral cross-sections.

\section*{Declaration of Generative AI in Scientific Writing}
During the preparation and development of this work, the author utilized generative AI assistants (Claude by Anthropic and Gemini by Google) for code refactoring, test generation, mathematical invariant verification, grammar refinement, and editorial structuring. The iterative review rounds conducted with AI models served as rigorous internal design audits and red-teaming exercises; they do not substitute for formal journal peer review. The author independently verified, executed, and validated all algorithms, mathematical proofs, test suites, and empirical benchmarks, and takes full responsibility for the factual and mathematical veracity of this publication.

\section*{Conflict of Interest}
The author develops and distributes educational dynamic geometry software, including packaged distributions of GeoSolver on the Microsoft Store. The core mathematical algorithms, classification engine, and provenance tracking system described in this paper are released as free and open-source software under the GNU General Public License v3 (GPLv3) without commercial restrictions.

\section*{Author Biography}
\textbf{Bao-Trung Nguyen} received his degree in Pharmacy from the University of Medicine and Pharmacy at Ho Chi Minh City in 2021. With over five years of professional experience in pharmaceutical research and development (R\&D), his current research interests include computational geometry, scientific software engineering, dynamic geometry systems, and robust arithmetic algorithms. He currently conducts scientific research as an independent researcher in Ho Chi Minh City, Vietnam.

\begin{thebibliography}{99}
\bibitem{shewchuk1997}
J.~R. Shewchuk, Adaptive precision floating-point arithmetic and fast robust geometric predicates, \textit{Discrete \& Computational Geometry} 18~(3) (1997) 305--363. \url{https://doi.org/10.1007/PL00009321}

\bibitem{cgal2024}
The CGAL Project, \textit{CGAL User and Reference Manual}, 5.6 edition, CGAL Editorial Board, 2024.

\bibitem{schwartz2010}
E.~J. Schwartz, T.~Avgerinos, D.~Brumley, All you ever wanted to know about dynamic taint analysis and forward symbolic execution (but might have been afraid to ask), in: \textit{2010 IEEE Symposium on Security and Privacy (S\&P)}, 2010, pp. 317--331. \url{https://doi.org/10.1109/SP.2010.26}

\bibitem{ieee754}
IEEE Computer Society, \textit{IEEE Standard for Floating-Point Arithmetic}, IEEE Std 754-2019, 2019, pp. 1--84. \url{https://doi.org/10.1109/IEEESTD.2019.8766229}

\bibitem{wcag21}
World Wide Web Consortium (W3C), Web Content Accessibility Guidelines (WCAG) 2.1, W3C Recommendation, 2018.
\end{thebibliography}

\end{document}
